\section{Introduction}
\label{sec:introduction}

\subsection{Background and Motivation}
\label{sec:background}

Modern offline Handwritten Text Recognition (HTR) has matured on the widely used benchmarks. A compact convolutional--recurrent CTC network with careful preprocessing~\cite{retsinas2022htrbestpractices} reaches roughly $4.6\%$ character error rate (CER) on IAM~\cite{marti2002iam}, and Transformer-based encoders such as HTR-VT~\cite{li2024htrvt} and pretrained encoder--decoder systems such as TrOCR~\cite{li2023trocr} reach comparable or lower error at higher parameter cost. What is not solved is what happens \emph{inside} these recognizers. Downstream document-analysis tasks --- writer identification, style analysis, palaeographic dating --- have traditionally relied on hand-designed features~\cite{souibgui2022vlac}, and a natural alternative is to reuse the internal attention of a Transformer HTR system directly. This alternative requires the attention maps to be clean: peaks that shift with the recognized character, low background diffusion, no spatially arbitrary artifacts.

The obstacle to reusing ViT attention was documented by Darcet~\etal~\cite{darcet2024registers}. In supervised and self-supervised Vision Transformers, a small number of patch tokens develop anomalously high norms and absorb a disproportionate share of the attention mass. These \emph{attention sinks} produce visible artifacts in attention maps and hurt any downstream task that reads those maps. Their proposed fix is minimally invasive: prepend a small number of learnable \emph{register tokens} to the patch sequence, and the sinks migrate onto the registers, leaving patch attention interpretable. On classification ViTs this fix is clean --- it improves downstream detection and segmentation while leaving top-1 accuracy unchanged --- but the mechanism has not been evaluated for CTC-based sequence-labelling HTR, where the queries that read the attention are the patch tokens themselves and no CLS-style aggregate query exists.

This project sits at that gap. The concrete deliverable is a modular HTR pipeline that trains a CNN-stem Vision Transformer with configurable register tokens (ViT-RGTS), keeps the CNN--RNN--CTC baseline of Retsinas~\etal\ as the reference point, and adds pretrained lanes (TorchVision ViT-B/16, TrOCR-base) for a fair architectural comparison. The concrete question is whether register tokens, evaluated at $R \in \{0, 2, 4, 8, 16\}$ on IAM and READ2016, change (a) the character/word error rate and (b) the quality of the attention maps produced by the encoder, judged both qualitatively in the style of Fig.~5 of \emph{Beyond Memorization}~\cite{nikolaidou2024beyondmem} and quantitatively through attention-concentration metrics.

\subsection{Problem Statement and Scope}
\label{sec:problem}

The project builds a Vision Transformer HTR system whose attention maps are fine-grained enough to support downstream writer-analysis tasks. Register tokens are the specific mechanism under test, motivated by the classification-ViT results of Darcet~\etal~\cite{darcet2024registers}. Success is measured jointly on character/word error rate and on attention quality; the writer-identification pipeline itself is out of scope for a 10~ECTS project budget and is sketched as future work.

The contributions of this report can be summarized in five sentences: I show that register tokens are CER-neutral within seed noise on both IAM and READ2016, matching the observation of Darcet~\etal\ for classification ViTs; I show that they reduce attention entropy, improve spatial ordering, and recruit more heads into the left-to-right reading pattern, while a residual sink survives on patch tokens; I reproduce the per-character attention visualization of Nikolaidou~\etal\ end-to-end on the register-added ViT; I identify the CNN stem (with strict left-to-right token order) as a structural requirement for CTC ViTs, evidenced by two independent architectural configurations that collapse to nearly $73\%$ CER; and I release the complete training and evaluation pipeline together with the run artifacts that produced these numbers, so that the writer-identification step suggested by the problem statement can be built directly on top.

The remainder of the report is organized as follows. Section~\ref{sec:related_work} surveys the HTR, register-token, and attention-interpretability literature and positions this work within it. Section~\ref{sec:methodology} describes the pipeline, datasets, and models. Section~\ref{sec:experiments} lays out the evaluation protocol. Section~\ref{sec:results} presents the recognition results, reproduces the \emph{Beyond Memorization} visualization on the trained model, and critically assesses whether the Darcet~\etal\ claim holds in the CTC HTR setting. Section~\ref{sec:discussion} interprets the findings, and Section~\ref{sec:conclusion} closes with the writer-identification extension.
